% Generated draft: requires mathematical and source review.
\documentclass[11pt]{article}
\usepackage[margin=25mm]{geometry}
\usepackage{fontspec}
\setmainfont{Latin Modern Roman}
\usepackage{amsmath,amssymb,mathtools,unicode-math}
\setmathfont{Latin Modern Math}
\usepackage{enumitem,needspace,xurl,hyperref}
\setlength{\parindent}{0pt}
\setlength{\parskip}{5pt}
\setlength{\emergencystretch}{3em}
\newcommand{\N}{\mathbb{N}}
\newcommand{\Z}{\mathbb{Z}}
\newcommand{\Q}{\mathbb{Q}}
\newcommand{\R}{\mathbb{R}}
\newcommand{\C}{\mathbb{C}}
\newcommand{\F}{\mathbb{F}}
\newcommand{\A}{\mathbb{A}}
\newcommand{\PP}{\mathbb{P}}
\newcommand{\E}{\mathbb{E}}
\newcommand{\eps}{\varepsilon}
\newcommand{\dd}{\,\mathrm{d}}
\begin{document}
\section*{UnsolvedMath ID: 11000128}
\textbf{Source problem code:} catalogue entry
\subsection{Definitions and mathematical statement}
Problem 4.2 asks for an analogue of the omitted Theorem 4.1 for topological manifolds, and for its relationship to surgery and pseudo-isotopy theory. The theorem itself, its manifold class, stabilization and conclusion are not supplied, so no specific topological analogue can be asserted from this record.

The subsequent background concerns a compact $d$-manifold $X$ and the pseudo-isotopy spaces
\[
 \operatorname{Diff}(X\times I\ \operatorname{rel}
 (\partial X\times I\cup X\times\{0\})),\qquad
 \operatorname{Top}(X\times I\ \operatorname{rel}
 (\partial X\times I\cup X\times\{0\})),
\]
where $I=[0,1]$ and the displayed subset is fixed pointwise. The top face is not added to the fixed subset. The smooth version requires a smooth structure on $X$. The source says Waldhausen's $A(X)$ explicitly contains information about both spaces in dimension-dependent ranges ([20], [21]), and less explicitly about $\operatorname{Diff}(X)$ and $\operatorname{Top}(X)$ ([22]); those ranges and comparison maps are omitted.

For a suitable based connected $X$, the usual spectral description is the algebraic $K$-theory of the suspension ring spectrum $\Sigma^\infty(\Omega X)_+$. The source describes its underlying infinite-loop object as the ``ring'' $\Omega^\infty S^\infty(\Omega X_+)$. Spectrum-versus-space and basepoint conventions must be fixed. It further describes topological cyclic homology $TC(X)$, introduced in [1], as related to cyclic homology, gives the cyclotomic trace $A(X)\to TC(X)$, and notes independent ring/linear-category applications ([9]).

The supplied formula (6), for a prime $p$, is
\[
 TC(\operatorname{pt})^\wedge_p\simeq
 \bigl(\Omega^\infty\mathbb S\times
 \Omega^{\infty-1}\mathbb{CP}_{-1}^\infty\bigr)^\wedge_p.
\]
Here the notation is interpreted in the customary stable sense: $\mathbb S$ is the sphere spectrum, and $\mathbb{CP}_{-1}^\infty$ conventionally denotes the Thom spectrum of the negative tautological complex line bundle. The source's exact completion and delooping conventions are absent. This formula and the free-loop-space description it mentions are contextual supplied assertions, not a replacement for the missing Theorem 4.1.
\subsection{Short English statement}
Find the topological-manifold counterpart of the omitted theorem and connect it with surgery and pseudo-isotopy theory. The supplied algebraic-$K$-theory and cyclic-homology background does not determine the missing target theorem.
\subsection{Sources}
\begin{itemize}
\item Local registry: \url{data/unsolvedmath/problems.json}, permanent UnsolvedMath ID $11000128$, exact source code \texttt{catalogue entry}.
\item Farb, \emph{Problems on Mapping Class Groups and Related Topics}, Problem 4.2, PDF page 105: \url{https://www.math.uchicago.edu/~farb/papers/mcgbook.pdf}. Supplied access date: 2026-07-29; extraction: pdf-font-structure; repeated local numbering disambiguated by page among 329 headings.
\item Source attributions and labels: Waldhausen's $A(X)$; [20], [21], [22] for diffeomorphism/homeomorphism comparisons; [1] for $TC$; [9] for further applications; equation (6). Full references and Theorem 4.1 are absent.
\item The registry attributes a smooth high-dimensional analogue to Galatius--Randal-Williams and says no topological solution was located, without an exact theorem or paper establishing the match to Theorem 4.1.
\end{itemize}
\end{document}
